% Options for packages loaded elsewhere
\PassOptionsToPackage{unicode}{hyperref}
\PassOptionsToPackage{hyphens}{url}
%
\documentclass[
]{article}
\usepackage{amsmath,amssymb}
\usepackage{iftex}
\ifPDFTeX
  \usepackage[T1]{fontenc}
  \usepackage[utf8]{inputenc}
  \usepackage{textcomp} % provide euro and other symbols
\else % if luatex or xetex
  \usepackage{unicode-math} % this also loads fontspec
  \defaultfontfeatures{Scale=MatchLowercase}
  \defaultfontfeatures[\rmfamily]{Ligatures=TeX,Scale=1}
\fi
\usepackage{lmodern}
\ifPDFTeX\else
  % xetex/luatex font selection
  \setmainfont[]{DejaVu Sans}
  \setmonofont[]{DejaVu Sans Mono}
\fi
% Use upquote if available, for straight quotes in verbatim environments
\IfFileExists{upquote.sty}{\usepackage{upquote}}{}
\IfFileExists{microtype.sty}{% use microtype if available
  \usepackage[]{microtype}
  \UseMicrotypeSet[protrusion]{basicmath} % disable protrusion for tt fonts
}{}
\makeatletter
\@ifundefined{KOMAClassName}{% if non-KOMA class
  \IfFileExists{parskip.sty}{%
    \usepackage{parskip}
  }{% else
    \setlength{\parindent}{0pt}
    \setlength{\parskip}{6pt plus 2pt minus 1pt}}
}{% if KOMA class
  \KOMAoptions{parskip=half}}
\makeatother
\usepackage{xcolor}
\usepackage[margin=24mm]{geometry}
\usepackage{color}
\usepackage{fancyvrb}
\newcommand{\VerbBar}{|}
\newcommand{\VERB}{\Verb[commandchars=\\\{\}]}
\DefineVerbatimEnvironment{Highlighting}{Verbatim}{commandchars=\\\{\}}
% Add ',fontsize=\small' for more characters per line
\newenvironment{Shaded}{}{}
\newcommand{\AlertTok}[1]{\textcolor[rgb]{1.00,0.00,0.00}{\textbf{#1}}}
\newcommand{\AnnotationTok}[1]{\textcolor[rgb]{0.38,0.63,0.69}{\textbf{\textit{#1}}}}
\newcommand{\AttributeTok}[1]{\textcolor[rgb]{0.49,0.56,0.16}{#1}}
\newcommand{\BaseNTok}[1]{\textcolor[rgb]{0.25,0.63,0.44}{#1}}
\newcommand{\BuiltInTok}[1]{\textcolor[rgb]{0.00,0.50,0.00}{#1}}
\newcommand{\CharTok}[1]{\textcolor[rgb]{0.25,0.44,0.63}{#1}}
\newcommand{\CommentTok}[1]{\textcolor[rgb]{0.38,0.63,0.69}{\textit{#1}}}
\newcommand{\CommentVarTok}[1]{\textcolor[rgb]{0.38,0.63,0.69}{\textbf{\textit{#1}}}}
\newcommand{\ConstantTok}[1]{\textcolor[rgb]{0.53,0.00,0.00}{#1}}
\newcommand{\ControlFlowTok}[1]{\textcolor[rgb]{0.00,0.44,0.13}{\textbf{#1}}}
\newcommand{\DataTypeTok}[1]{\textcolor[rgb]{0.56,0.13,0.00}{#1}}
\newcommand{\DecValTok}[1]{\textcolor[rgb]{0.25,0.63,0.44}{#1}}
\newcommand{\DocumentationTok}[1]{\textcolor[rgb]{0.73,0.13,0.13}{\textit{#1}}}
\newcommand{\ErrorTok}[1]{\textcolor[rgb]{1.00,0.00,0.00}{\textbf{#1}}}
\newcommand{\ExtensionTok}[1]{#1}
\newcommand{\FloatTok}[1]{\textcolor[rgb]{0.25,0.63,0.44}{#1}}
\newcommand{\FunctionTok}[1]{\textcolor[rgb]{0.02,0.16,0.49}{#1}}
\newcommand{\ImportTok}[1]{\textcolor[rgb]{0.00,0.50,0.00}{\textbf{#1}}}
\newcommand{\InformationTok}[1]{\textcolor[rgb]{0.38,0.63,0.69}{\textbf{\textit{#1}}}}
\newcommand{\KeywordTok}[1]{\textcolor[rgb]{0.00,0.44,0.13}{\textbf{#1}}}
\newcommand{\NormalTok}[1]{#1}
\newcommand{\OperatorTok}[1]{\textcolor[rgb]{0.40,0.40,0.40}{#1}}
\newcommand{\OtherTok}[1]{\textcolor[rgb]{0.00,0.44,0.13}{#1}}
\newcommand{\PreprocessorTok}[1]{\textcolor[rgb]{0.74,0.48,0.00}{#1}}
\newcommand{\RegionMarkerTok}[1]{#1}
\newcommand{\SpecialCharTok}[1]{\textcolor[rgb]{0.25,0.44,0.63}{#1}}
\newcommand{\SpecialStringTok}[1]{\textcolor[rgb]{0.73,0.40,0.53}{#1}}
\newcommand{\StringTok}[1]{\textcolor[rgb]{0.25,0.44,0.63}{#1}}
\newcommand{\VariableTok}[1]{\textcolor[rgb]{0.10,0.09,0.49}{#1}}
\newcommand{\VerbatimStringTok}[1]{\textcolor[rgb]{0.25,0.44,0.63}{#1}}
\newcommand{\WarningTok}[1]{\textcolor[rgb]{0.38,0.63,0.69}{\textbf{\textit{#1}}}}
\setlength{\emergencystretch}{3em} % prevent overfull lines
\providecommand{\tightlist}{%
  \setlength{\itemsep}{0pt}\setlength{\parskip}{0pt}}
\setcounter{secnumdepth}{-\maxdimen} % remove section numbering
\ifLuaTeX
  \usepackage{selnolig}  % disable illegal ligatures
\fi
\usepackage{bookmark}
\IfFileExists{xurl.sty}{\usepackage{xurl}}{} % add URL line breaks if available
\urlstyle{same}
\hypersetup{
  hidelinks,
  pdfcreator={LaTeX via pandoc}}

\author{}
\date{}

\begin{document}

\section{ClientOpsCRM 2.0.0 Open-Source Release
Guide}\label{clientopscrm-200-open-source-release-guide}

\subsection{What the source AIO
contains}\label{what-the-source-aio-contains}

The source AIO retains the complete ClientOps Relay 1.2.1 GitHub release
and adds the CRM source, build system, tests, Relay bridge,
control/planner material, research corpus, CRM documentation and
verification evidence. The inherited Apache-2.0 license is preserved.

\subsection{Rebuild from source}\label{rebuild-from-source}

\begin{Shaded}
\begin{Highlighting}[]
\FunctionTok{make} \AttributeTok{{-}C}\NormalTok{ crm clean release debug sanitize}
\ExtensionTok{python3}\NormalTok{ crm/tests/test\_crm.py}
\VariableTok{CRM\_BIN}\OperatorTok{=}\NormalTok{crm/build/clientops{-}crm{-}sanitize }\ExtensionTok{python3}\NormalTok{ crm/tests/test\_crm.py}
\end{Highlighting}
\end{Shaded}

The release pack includes a Linux x86-64 build for convenience, but
source is canonical. Other systems should rebuild with a C++17 compiler
and SQLite development library.

\subsection{Recommended repository
root}\label{recommended-repository-root}

Publish the contents of the source AIO as the repository root. Do not
commit generated local state (\texttt{runtime/}), \texttt{.env},
dependency directories, compiler scratch or private deployment secrets.
The shipped \texttt{.gitignore} covers CRM build/runtime paths in
addition to Relay exclusions.

\subsection{Primary entry points}\label{primary-entry-points}

\begin{itemize}
\tightlist
\item
  \texttt{README.md} --- combined project start.
\item
  \texttt{clientops} --- combined operator launcher.
\item
  \texttt{crm/src/clientops\_crm.cpp} --- CRM implementation.
\item
  \texttt{crm/Makefile} --- CRM build gates.
\item
  \texttt{crm/tests/test\_crm.py} --- real integration/adversarial test.
\item
  \texttt{crm/scripts/import-relay.sh} --- live Relay-to-CRM bridge.
\item
  \texttt{policy-engine/} --- retained Relay C11 deterministic policy
  engine.
\item
  \texttt{database/} --- retained Relay PostgreSQL schema/migrations.
\item
  \texttt{workflows/} --- retained Relay n8n workflows.
\item
  \texttt{research/} --- 143-source research corpus and claim ledger.
\item
  \texttt{docs/crm/} --- CRM documentation.
\item
  \texttt{evidence/crm/} --- final machine verification evidence.
\end{itemize}

\subsection{Release discipline}\label{release-discipline}

\begin{enumerate}
\def\labelenumi{\arabic{enumi}.}
\tightlist
\item
  Rebuild from a clean tree.
\item
  Run the CRM release and sanitizer suites.
\item
  Run Relay-compatible gates available on the target toolchain.
\item
  Run \texttt{./clientops\ crm\ -\/-json\ doctor} on a fresh database.
\item
  Produce a deterministic release package.
\item
  Verify ZIP integrity.
\item
  Verify no empty files, secrets, runtime DBs or dependency directories
  are shipped.
\item
  Publish SHA-256 manifest alongside the release asset.
\end{enumerate}

\subsection{Versioning}\label{versioning}

The combined product identifies itself as
\texttt{ClientOpsCRM\ 2.0.0\ +\ ClientOps\ Relay\ 1.2.1}.
Relay\textquotesingle s \texttt{package.json} stays at 1.2.1 so its own
release/build tooling continues to describe the inherited subsystem
correctly; \texttt{CLIENTOPSCRM\_VERSION} carries the combined release
identity.

\end{document}
